\chapter{Dataset construction and shared representation}\label{sec:dataset}

This chapter defines the \emph{data} artifact produced by the project. Chapter~\ref{sec:system} described the applications that create and consume data; here the focus is the common representation that lets two different construction paths, controlled SUMO simulation and conversion of observed GPS trajectories, feed the same training and web-evaluation pipeline.

\section{Purpose and scope}
The target artifact is a route-aware, time-indexed traffic graph for ETA-oriented learning. It must preserve enough road topology, vehicle state, route context, and future-arrival information for a downstream model to learn travel-time structure, while remaining independent of how the source traffic was produced. This separation is the main invariant of the dataset layer: simulation-derived examples and trajectory-derived examples may differ in observability and noise, but after construction they expose one contract to model code and to the \WebApp{}.

The representation is also deliberately richer than a table of final travel times. It stores static network structure, dynamic vehicle nodes, route membership, road-level traffic descriptors, and labels at sampled times. As a result, the same dataset can support vehicle-level ETA prediction and related analyses such as road load, region-level congestion, or ablation studies over route features.

\section{Two construction paths, one output contract}
\paragraph{Simulation-derived construction.}
The simulation path starts from a SUMO scenario controlled by the desktop tool and TraCI~\cite{sumo2011,traci2018}. This branch has strong observability: the network, planned routes, vehicle identities, positions, speeds, and simulator time are all available during execution. The desktop configuration defines zones, landmarks, vehicle classes, schedules, sampling cadence, output options, and optional perturbations such as untracked random trips. The operator-facing steps are documented in Section~\ref{sec:system:desktop:ui}; Chapter~4 uses those screens only as evidence and does not repeat them.

During a run, the generator samples simulator state at the selected interval and writes a sequence of graph snapshots. Each snapshot records the active vehicles and the road segments carrying traffic or demand at that time. Because the simulation clock is controlled, this branch is useful for reproducible benchmark settings and for experiments where the amount of traffic or the vehicle-type mixture must be known in advance.

\paragraph{Trajectory-derived construction.}
The trajectory path starts from GPS-like rows and a SUMO-compatible road network prepared from OSM~\cite{osm}. It is less controlled than simulation: traces may include sparse samples, repeated start/end points, unrealistic jumps, and projection errors. The conversion workflow therefore performs normalization, route matching, invalid-segment handling, and projection of GPS points onto inferred SUMO routes before emitting graph snapshots. The Porto taxi data~\cite{porto2015pkdd} is used as the representative bundled example, but the contract is intentionally general.

For successful conversion, imported CSV files must provide one trajectory polyline per row in a parseable list form (the converter scans for a cell like \texttt{[[...], ...]}), with each point encoded as a coordinate pair and at least two points per trip. A \texttt{TIMESTAMP} column is optional for inspecting individual routes but required by the batch step-conversion path that emits sampled snapshots and labels; it anchors each trajectory on a global time axis. Additional source columns may exist, but conversion consumes only the trajectory geometry and, when present, the timestamp. External datasets should therefore be normalized to this minimal contract before import. Section~\ref{sec:dataset:csv-contract} gives the concrete CSV shape.

\paragraph{Convergence.}
Both branches end in the same exported layout: static network material, time-step snapshots, and labels. This is why the web application is independent of the construction path (Figure~\ref{fig:web_ui_porto}) and why evaluation results can compare simulation-backed and trajectory-backed data without path-specific adapters.

The generated graph-representation datasets for the 3-zone simulation network and the Porto taxi conversion example are publicly available in a \href{https://drive.google.com/drive/folders/1baEONbv8WyDVanF_qvLCq1UJ4oT5N7LE?usp=sharing}{shared dataset folder}. They are distributed under an attribution license: users may reuse the datasets, but should cite the author and this work when doing so. The desktop dataset generator and trajectory-conversion tool are available in the public repository \href{https://github.com/Turgibot/Multi-Variant-Simulated-Traffic-Dataset-Creator-and-Model-Tester}{Turgibot/Multi-Variant-Simulated-Traffic-Dataset-Creator-and-Model-Tester}, so users can reproduce the bundled examples or construct graph datasets from their own networks and trajectory sources.

\section{Trajectory CSV input contract}\label{sec:dataset:csv-contract}
The trajectory-conversion branch accepts CSV files whose rows describe complete trips. The converter does not require the full Porto schema; it only needs the geometry and, for batch snapshot generation, a start timestamp. Therefore, other trajectory datasets can be used after being normalized to the following minimal structure.

\begin{table}[!htbp]
\centering
\scriptsize
\begin{tabular}{p{0.22\linewidth}p{0.20\linewidth}p{0.46\linewidth}}
\toprule
\textbf{CSV field} & \textbf{Required} & \textbf{Meaning} \\
\midrule
\texttt{TIMESTAMP} & Required for batch export & Integer start time for the trip. It is used to place the converted route on the global time axis before sampling snapshots. \\
Polyline cell & Required & Any cell in the row containing a parseable list of coordinate pairs, for example \texttt{[[-8.61,41.15],[-8.60,41.16]]}. The converter scans cells for this shape. \\
Other columns & Optional & Source identifiers, taxi metadata, missing-data flags, or dataset-specific fields. They may remain in the CSV, but they are not consumed by conversion. \\
\bottomrule
\end{tabular}
\caption{Minimal trajectory CSV contract for conversion.}
\label{tab:trajectory_csv_contract}
\end{table}

A compact valid row is shown below. The coordinate order is the one expected by the converter: longitude followed by latitude in each point. The example contains only the fields that matter for conversion; real source files may contain additional columns.

\begin{lstlisting}[caption={Minimal trajectory CSV example.},label={lst:trajectory-csv-example}]
TRIP_ID,TIMESTAMP,POLYLINE
1,1372636858,"[[-8.6186,41.1413],[-8.6179,41.1420],[-8.6168,41.1431]]"
\end{lstlisting}

A trajectory is therefore interpreted as an ordered polyline
\[
P = (p_1,\ldots,p_m), \quad p_i=(\mathrm{lon}_i,\mathrm{lat}_i),
\]
with an associated start time. During conversion, repeated stationary endpoints may be trimmed, unrealistic jumps may split $P$ into smaller valid segments, and each remaining segment is matched to the SUMO network.

\section{Exported graph artifact}\label{sec:shared-representation}

% Shared graph layers (Chapter: dataset and representation)
\begin{figure}[!htbp]
\centering
\includegraphics[width=0.5\linewidth,keepaspectratio]{dynamic_static_nodes_and_edges.png}
\caption{Static vs.\ dynamic graph layers. Static layer: black nodes A--D are junctions; grey edges are roads. Dynamic layer: green nodes are vehicles; dashed blue edges connect vehicles to other vehicles and to junctions.}
\label{fig:graph_schema}
\end{figure}


Figure~\ref{fig:graph_schema} summarizes the layer split. The static layer contains junction nodes and directed road edges. The dynamic layer contains vehicle nodes and temporary relations created by vehicle movement and proximity. Figure~\ref{fig:traffic_graph_snapshot} shows the same idea on a concrete morning traffic snapshot: the road graph forms the stable substrate, active vehicles are placed along the network, and route/dynamic relations connect the current traffic state to downstream ETA labels. In the implementation, the export is organized as:

\begin{figure}[!htbp]
\centering
\includegraphics[width=\linewidth,keepaspectratio]{graph_representation_of_traffic.png}
\caption{Graph representation of a morning traffic snapshot. Blue nodes are static junctions, green nodes are dynamic vehicle nodes, grey lines are static road edges, and orange lines are dynamic edges that encode route-aware traffic relations.}
\label{fig:traffic_graph_snapshot}
\end{figure}

\begin{itemize}
    \item \textbf{Static network file} (\texttt{static.json} or \texttt{static.json.gz}): junctions and static road-edge attributes such as identifiers, geometry-derived structure, lengths, and stable network metadata.
    \item \textbf{Snapshot files} (\texttt{snapshots/step\_*.json} or compressed equivalents): a sampled time step with active vehicle nodes, dynamic road-edge state, and dynamic vehicle relations.
    \item \textbf{Label files} (\texttt{labels/label\_*.json} or compressed equivalents): the supervised target for active vehicles at the same timestamp, primarily ETA in seconds.
\end{itemize}

The static file is shared by all time steps in a generated dataset. Snapshot files carry evolving state: vehicle position, speed, acceleration, current edge, remaining route, route length left, road demand, average speed, density, and vehicle-to-road or vehicle-to-vehicle context as available. Label files align with snapshots by timestamp and store, for each vehicle with a known destination time, the remaining time to arrival. This keeps features and targets synchronized without baking the target into the feature graph.

\section{Entity feature schema}
The exported representation uses four graph entity families: junctions, road edges, vehicle nodes, and dynamic edges. Static junctions and static road-edge fields are written once in \texttt{static.json}; vehicle nodes, dynamic road-edge fields, and dynamic edges are written per snapshot; ETA labels are written in matching label files.

\begin{table}[!htbp]
\centering
\scriptsize
\begin{tabular}{p{0.24\linewidth}p{0.18\linewidth}p{0.46\linewidth}}
\toprule
\textbf{Junction feature} & \textbf{Type} & \textbf{Description} \\
\midrule
\texttt{id} & string & Stable junction identifier from the SUMO network. \\
\texttt{x} & number & Junction x-position in SUMO/map coordinates. \\
\texttt{y} & number & Junction y-position in SUMO/map coordinates. \\
\texttt{type} & string & SUMO junction type, defaulting to \texttt{priority} when unavailable. \\
\texttt{zone} & string & Optional zone label used by scenario configuration. \\
\texttt{incoming} & list[string] & Incoming road-edge identifiers. \\
\texttt{outgoing} & list[string] & Outgoing road-edge identifiers. \\
\bottomrule
\end{tabular}
\caption{Static junction-node features in \texttt{static.json}.}
\label{tab:junction_features}
\end{table}

\begin{table}[!htbp]
\centering
\scriptsize
\begin{tabular}{p{0.24\linewidth}p{0.18\linewidth}p{0.46\linewidth}}
\toprule
\textbf{Road feature} & \textbf{Type} & \textbf{Description} \\
\midrule
\texttt{id} & string & Stable road-edge identifier. Lane suffixes are collapsed to a base edge ID. \\
\texttt{from} & string & Junction identifier at the start of the directed edge. \\
\texttt{to} & string & Junction identifier at the end of the directed edge. \\
\texttt{edge\_type} & integer & Static road edges use \texttt{0}; dynamic relation types are listed separately. \\
\texttt{speed} & number & SUMO speed limit or representative free-flow speed. \\
\texttt{length} & number & Edge length in metres/map units as supplied by SUMO. \\
\texttt{num\_lanes} & integer & Number of lanes represented by the collapsed road edge. \\
\texttt{zone} & string & Optional zone assignment. \\
\texttt{vehicles\_on\_road} & list[string] & Dynamic list of active vehicle IDs currently on the edge. \\
\texttt{edge\_demand} & number & Dynamic downstream demand contribution from active routes. \\
\texttt{avg\_speed} & number & Dynamic average speed of vehicles currently on the edge. \\
\texttt{density} & number & Dynamic occupancy estimate, computed from vehicles, length, and lane count. \\
\bottomrule
\end{tabular}
\caption{Road-edge features. Static fields appear in \texttt{static.json}; dynamic traffic fields appear in each snapshot under \texttt{road\_edges\_dynamic}.}
\label{tab:road_features}
\end{table}

\begin{table}[!htbp]
\centering
\scriptsize
\begin{tabular}{p{0.24\linewidth}p{0.18\linewidth}p{0.46\linewidth}}
\toprule
\textbf{Vehicle feature} & \textbf{Type} & \textbf{Description} \\
\midrule
\texttt{id} & string & Vehicle/node identifier. \\
\texttt{vehicle\_type} & string & Vehicle class, e.g. passenger, bus, or truck. \\
\texttt{length} & number & Vehicle length when known. \\
\texttt{width} & number & Vehicle width when known. \\
\texttt{height} & number & Vehicle height when known. \\
\texttt{speed} & number & Current vehicle speed. \\
\texttt{acceleration} & number & Current vehicle acceleration. \\
\texttt{current\_x} & number & Current projected x-coordinate. \\
\texttt{current\_y} & number & Current projected y-coordinate. \\
\texttt{current\_zone} & string & Zone currently occupied, when available. \\
\texttt{current\_edge} & string & Road edge currently occupied. \\
\texttt{current\_position} & number & Position along the current edge. \\
\texttt{origin\_name} & string & Human-readable origin label, when available. \\
\texttt{origin\_zone} & string & Zone assigned to the trip origin. \\
\texttt{origin\_edge} & string & Road edge containing the origin point. \\
\texttt{origin\_position} & number & Position along the origin edge. \\
\texttt{origin\_x} & number & Origin x-coordinate. \\
\texttt{origin\_y} & number & Origin y-coordinate. \\
\texttt{origin\_start\_sec} & integer & Trip start time in seconds on the dataset time axis. \\
\texttt{route} & list[string] & Full planned or inferred edge route. \\
\texttt{route\_length} & number & Total route length. \\
\texttt{route\_left} & list[string] & Remaining route after the current edge. \\
\texttt{route\_length\_left} & number & Remaining distance estimate to destination. \\
\texttt{destination\_name} & string & Human-readable destination label, when available. \\
\texttt{destination\_edge} & string & Road edge containing the destination point. \\
\texttt{destination\_position} & number & Position along the destination edge. \\
\texttt{destination\_x} & number & Destination x-coordinate. \\
\texttt{destination\_y} & number & Destination y-coordinate. \\
\bottomrule
\end{tabular}
\caption{Vehicle-node features in snapshot \texttt{nodes}.}
\label{tab:vehicle_features}
\end{table}

\begin{table}[!htbp]
\centering
\scriptsize
\begin{tabular}{p{0.24\linewidth}p{0.18\linewidth}p{0.46\linewidth}}
\toprule
\textbf{Dynamic-edge feature} & \textbf{Type} & \textbf{Description} \\
\midrule
\texttt{id} & string & Generated relation identifier. \\
\texttt{from} & string & Source node identifier. \\
\texttt{to} & string & Target node identifier. \\
\texttt{edge\_type} & integer & \texttt{0}: static road edge; \texttt{1}: junction-to-first-vehicle; \texttt{2}: vehicle-to-vehicle; \texttt{3}: last-vehicle-to-junction. \\
\bottomrule
\end{tabular}
\caption{Dynamic relation-edge features in snapshot \texttt{dynamic\_edges}.}
\label{tab:dynamic_edge_features}
\end{table}

\begin{table}[!htbp]
\centering
\scriptsize
\begin{tabular}{p{0.24\linewidth}p{0.18\linewidth}p{0.46\linewidth}}
\toprule
\textbf{Label field} & \textbf{Type} & \textbf{Description} \\
\midrule
\texttt{timestamp} & integer & Snapshot time that the labels align to. \\
\texttt{labels[].id} & string & Vehicle identifier. \\
\texttt{labels[].eta} & number & Remaining time to destination in seconds, clipped at zero after arrival. \\
\bottomrule
\end{tabular}
\caption{ETA label-file schema.}
\label{tab:label_features}
\end{table}

\FloatBarrier

\section{Compact output example}
The implementation writes static data, snapshots, and labels as separate files. Listing~\ref{lst:dataset-json-example} shows a shortened, aligned excerpt with one junction, one road edge, one vehicle, one dynamic edge, and one ETA label. Values are illustrative; the structure mirrors the generated artifact.

\begin{lstlisting}[caption={Compact dataset JSON excerpt across \texttt{static.json}, one snapshot, and one label file.},label={lst:dataset-json-example}]
{
  "static.json": {
    "junctions": [
      {
        "id": "J1",
        "x": 10.0,
        "y": 20.0,
        "type": "priority",
        "zone": "A",
        "incoming": [],
        "outgoing": ["E1"]
      }
    ],
    "road_edges": [
      {
        "id": "E1",
        "from": "J1",
        "to": "J2",
        "edge_type": 0,
        "speed": 13.9,
        "length": 120.0,
        "num_lanes": 1,
        "zone": "A"
      }
    ]
  },
  "snapshots/step_000000001030.json": {
    "step": 1030,
    "nodes": [
      {
        "id": "veh_1",
        "vehicle_type": "passenger",
        "speed": 8.4,
        "acceleration": 0.2,
        "current_x": 42.0,
        "current_y": 77.0,
        "current_zone": "A",
        "current_edge": "E1",
        "current_position": 35.0,
        "origin_edge": "E1",
        "origin_position": 0.0,
        "destination_edge": "E9",
        "destination_position": 70.0,
        "route": ["E1", "E5", "E9"],
        "route_length": 420.0,
        "route_left": ["E5", "E9"],
        "route_length_left": 385.0
      }
    ],
    "road_edges_dynamic": [
      {
        "id": "E1",
        "vehicles_on_road": ["veh_1"],
        "edge_demand": 0.8,
        "avg_speed": 8.4,
        "density": 0.0083
      }
    ],
    "dynamic_edges": [
      {
        "id": "dyn_000000001",
        "from": "J1",
        "to": "veh_1",
        "edge_type": 1
      }
    ]
  },
  "labels/label_000000001030.json": {
    "timestamp": 1030,
    "labels": [
      {
        "id": "veh_1",
        "eta": 420
      }
    ]
  }
}
\end{lstlisting}

\section{Temporal alignment and labels}
Time is the bridge between construction and learning. In the simulation branch, the source clock is the SUMO simulation time and snapshots are sampled at fixed intervals. In the trajectory branch, each trajectory row is first converted into a sequence of projected points with timestamps; the converter then aggregates these point updates into the same fixed-period snapshot schedule.

For ETA learning, a vehicle label at time $t$ is the remaining time until that vehicle's destination event, clipped at zero after arrival. This target can be generated directly in simulation because the run observes vehicle completion. For converted trajectories, the destination event is inferred from the last valid projected point in the matched trajectory segment. The important property is that both branches produce labels with the same meaning: \emph{given the graph state and route context at timestamp $t$, estimate remaining travel time for the vehicle}.

\section{Trajectory repair and route matching}
Trajectory conversion must turn noisy GPS polylines into routes on the SUMO network. The process has three roles: remove or split invalid observations, infer a plausible network path, and project samples onto that path so graph snapshots can be generated.

First, repeated start/end points can be trimmed so a trip has a single effective origin and destination. Unrealistic jumps can be treated as invalid segments; when such jumps occur, the trajectory is split and each valid segment is matched separately. The GUI diagnostics in Section~\ref{sec:system:desktop:ui} show these controls and their visual outputs without duplicating figures in this chapter.

Matching is not run on the full regional road graph. An axis-aligned bounding box $B$ is built around the GPS polyline $P$ (optionally padded), and $N_{\mathrm{bb}}=(V_{\mathrm{bb}},E_{\mathrm{bb}})$ is the subnetwork formed from all edges whose geometry intersects $B$ and their incident nodes. GPS polylines are then aligned to $N_{\mathrm{bb}}$ in two phases: a \emph{scoring pass} assigns every edge a coarse tier ($1$, $10$, or $1000$) from successive segment-to-edge comparisons, and a \emph{routing pass} searches for a low-cost path that respects those tiers.

For each polyline segment, edges whose heading differs from the segment by more than $\pi/9$ are discarded. Among survivors, edges are ranked by
\[
r_e=\lambda\delta_e+\frac{1}{2}\bigl(d(p_i,e)+d(p_{i+1},e)\bigr),
\]
where $\delta_e$ is the acute angle gap in radians, $d(p,e)$ is the distance from point $p$ to the centreline polyline of edge $e$, and $\lambda$ fixes the radian-to-map-unit trade-off. The best $k$ edges become primary candidates and the next $j$ become secondary candidates. A multi-start A$^\ast$ pass then searches between primary candidates near the origin and primary candidates near the destination. Algorithms~\ref{alg:route-edge-weights}--\ref{alg:route-astar} state this procedure formally.

% Route matching: abstract scoring + A* as specified for the paper (conceptual layer).

\begin{algorithm}[htbp]
\caption{Edge-weight assignment from a GPS polyline (scoring pass)}
\label{alg:route-edge-weights}
\begin{algorithmic}[1]
\Require Polyline $P=(p_1,\ldots,p_m)$; full road network $N=(V,E)$ with directed edges;
         integers $k,j \ge 1$; angular threshold $\theta_{\max}=\pi/9$; trade-off $\lambda>0$ (same length units as the map, e.g.\ $\lambda=5$).
\Ensure Each edge $e\in E_{\mathrm{bb}}$ has a tier $w(e)\in\{1,10,1000\}$ (lower is more attractive).
\State Build axis-aligned bounding box $B$ containing all points of $P$ (optional margin).
\State Let $N_{\mathrm{bb}}=(V_{\mathrm{bb}},E_{\mathrm{bb}})\subseteq N$ be the subnetwork of all edges in $E$ whose centreline intersects $B$, together with their endpoint nodes in~$V$.
\State Initialise $w(e)\gets 1000$ for all $e\in E_{\mathrm{bb}}$.
\For{each consecutive segment $s_i=(p_i,p_{i+1})$, $i=1,\ldots,m-1$}
    \State $\phi_i \gets \mathrm{atan2}\bigl(p_{i+1}-p_i\bigr)$ \Comment{segment heading in the projected map}
    \State $\mathcal{C}\gets\emptyset$ \Comment{candidate edges for this segment}
    \For{each edge $e\in E_{\mathrm{bb}}$}
        \State $\psi_e \gets$ heading of $e$ from its centreline polyline (first to last shape point)
        \State $\delta_e \gets \min\bigl(|\phi_i-\psi_e|,\,2\pi-|\phi_i-\psi_e|\bigr)$
        \If{$\delta_e \le \theta_{\max}$}
            \State mark $e$ as \emph{potential} for $s_i$; add $e$ to $\mathcal{C}$
            \State $d_i \gets d(p_i,e)$; $d_{i+1} \gets d(p_{i+1},e)$
            \Comment{shortest Euclidean distance from each segment endpoint to the polyline of $e$}
            \State $r_e \gets \lambda\,\delta_e + \tfrac{1}{2}(d_i + d_{i+1})$
            \Comment{$\lambda$ scales the angle gap $\delta_e$ (radians) into map units so it is comparable to the mean endpoint--edge distances.}
        \EndIf
    \EndFor
    \State sort $\mathcal{C}$ by increasing $r_e$
    \State $\mathrm{Top}_i \gets$ first $k$ edges; $\mathrm{Sec}_i \gets$ next $j$ edges (if $|\mathcal{C}|<k{+}j$, take as many as exist)
    \For{$e\in \mathrm{Top}_i$}
        \State $w(e)\gets 1$
    \EndFor
    \For{$e\in \mathrm{Sec}_i$}
        \State $w(e)\gets \min\bigl(w(e),10\bigr)$
    \EndFor
\EndFor
\Comment{Edges that were never marked potential in any segment keep $w(e)=1000$. Each time an edge ranks in $\mathrm{Top}_i$, set $w(e)\gets 1$; secondary ranks use $\min$ with $10$ so a prior tier $1$ is not overwritten by a weaker segment.}
\end{algorithmic}
\end{algorithm}

\begin{algorithm}[htbp]
\caption{Route extraction by multi-start A$^\ast$}
\label{alg:route-astar}
\begin{algorithmic}[1]
\Require Restricted network $N_{\mathrm{bb}}$; tiers $w(e)$ on $E_{\mathrm{bb}}$; sets $\mathrm{Top}_1$ and $\mathrm{Top}_{m-1}$ from Algorithm~\ref{alg:route-edge-weights} (primaries for the first and last polyline segments); last point $p_m$.
\Ensure Shortest-cost route among $k$ multi-start runs.
\State Goal set $G \gets \mathrm{Top}_{m-1}$ \Comment{the $k$ primary edges for the last segment $(p_{m-1},p_m)$}
\For{each start edge $e_{\mathrm{start}}\in \mathrm{Top}_1$}
    \State Run A$^\ast$ from $e_{\mathrm{start}}$ with edge cost $c(e)=w(e)\cdot \mathrm{len}(e)$ and $h(v)=$ Euclidean distance from $v$ to $p_m$;
           \textbf{stop} the run as soon as the search reaches any state on an edge in $G$.
\EndFor
\State \Return the route with minimum total cost among these $k$ runs.
\end{algorithmic}
\end{algorithm}


After a route is inferred, GPS samples are projected onto the SUMO route. The converter records the current edge, position along the edge, speed, distance from previous point, route edges, and remaining route context. These projected states are what make trajectory-derived data compatible with simulation-derived snapshots.

\FloatBarrier

\section{Representation invariants}
The shared contract rests on a small set of invariants:
\begin{itemize}
    \item \textbf{One static substrate:} road and junction identifiers remain stable across all snapshots in a dataset.
    \item \textbf{Time-indexed dynamics:} vehicle nodes and road-edge dynamic attributes are sampled at explicit timestamps.
    \item \textbf{Route conditioning:} each tracked vehicle carries an origin, destination, planned or inferred route, and route-left information.
    \item \textbf{Separated labels:} ETA targets are stored alongside snapshots but not mixed into static or dynamic feature fields.
    \item \textbf{Path independence:} downstream consumers see the same graph schema whether the source was a simulator run or a converted trajectory file.
\end{itemize}

These invariants are the reason Chapter~\ref{sec:web} can treat both branches through the same evaluation client. The web application loads a compatible exported representation, computes or displays routes, invokes an ETA predictor, and stores the logged prediction and outcome evidence without needing to know whether the original data came from controlled SUMO generation or real-world GPS conversion.
